% Proposed narrative revision based on frozen commit
% 986149c07bc7b8be54d99e8a2e2a47327ec143b5.
% Uses bibliography keys already present in the frozen submission.
% This draft does not assert FPGA timing closure or universal superiority.
% Authors must verify all wording against the final selected implementation.
\section{Introduction}
Three-dimensional and multi-die FPGAs expand the design space of programmable interconnects, while placing new demands on the CAD tools used to evaluate them. Recent frameworks support flexible architecture generation, stacked-fabric placement, and multi-die-aware routing~\cite{lazagna,flatland,multidie}. These flows rely on router lookahead tables that summarize expected delay and congestion costs for different routing resources and destination contexts. Efficiently constructing these tables matters when an architecture or cost-model change invalidates an existing table. Compatible precomputed tables should be reused; our focus is the construction work that remains when such reuse is unavailable.

However, parallel routing does not automatically make lookahead construction parallel. In our locked VPR flow, an eight-core GEMM experiment spends 149.56 seconds constructing lookahead---approximately 45\% of its fixed-input process time. Two execution bottlenecks restrict concurrency: barriers between groups containing only a few samples, and fine-grained priority ordering within each search. Simply merging searches or relaxing their ordering is not a drop-in solution. Each sample retains its source and layer-relative recording context, and a returned delay must remain associated with the congestion representative selected under the relevant search rule. The challenge is therefore to expose parallel work without confusing removable scheduling constraints with the dependencies that define the output table.

Existing approaches address complementary parts of this challenge. MultiQueue-based routing parallelizes connection searches through carefully controlled priority relaxation~\cite{multiqueue}; it does not directly provide the sample-recording and reduction semantics of lookahead construction. Multi-die-aware lookahead improves how heuristic costs represent die crossings~\cite{multidie}, addressing the heuristic itself rather than the execution organization studied here. Parallel shortest-path methods already exploit buckets and frontier expansion to balance useful work against scheduling overhead~\cite{delta,gpu_sssp}. Applying these mechanisms to lookahead nevertheless requires a concrete treatment of source-relative table entries, coupled delay/congestion labels, same-bucket completion, and the native consumer interface. Table reuse and sparse-sampling lookahead variants provide additional, complementary ways to avoid or reduce construction work~\cite{vproptions}; they do not remove the need to parallelize the searches retained by a chosen construction policy.

We address this problem through a two-granularity decomposition that separates logical sample identity from execution order. STRICT exposes samples across group boundaries, uses private mutable workspaces, and restores native-order reduction under explicit wire-map equivalence conditions. ALT separately defines FP32 whole-pair bucket waves, expanding search-internal concurrency while preserving the association between each delay and congestion value; its differences from native selection are evaluated rather than hidden. A common task boundary connects these methods to the native builder, and a finite FPGA mapping investigates the corresponding ownership, memory, and retirement requirements. The measured benefits extend beyond isolated kernels: STRICT Pipeline achieves a $1.638\times$ Process speedup over Stock on the eight-core GEMM cohort, while ALT CUDA achieves $3.207\times$ native construction and $1.094\times$ Process speedups over optimized CPU Pipeline on the same-family LSTM tuning holdout. These are separate measured cohorts. Together with controlled scheduling and wave-work comparisons, they show how removing specific execution constraints can accelerate lookahead construction while keeping its semantic and system-level consequences explicit.
